%------------------------------------------------------------------------------
% Início do resumo
%------------------------------------------------------------------------------
Amaro, Felipe Pereira~\textbf{Inteligência artificial na escrita científica: um tutor de gênero automatizado para submissões a conferências brasileiras}. UNIRIO, 2026. \pageref{LastPage} páginas. Qualificação de Mestrado. Programa de Pós-Graduação em Informática, UNIRIO.
\vspace{10pt}
\begin{center}
    \textbf{RESUMO}
    \vspace{10pt}
\end{center}

Este trabalho investiga o uso de modelos de linguagem de grande escala (LLMs) como instrumento de apoio à escrita científica no contexto de conferências brasileiras. O crescimento da produção científica intensificou limitações estruturais do processo de revisão por pares, como a sobrecarga de revisores, a heterogeneidade dos pareceres e a escassez de retorno estruturado a autores em início de carreira. Uma consequência particularmente onerosa para novos pesquisadores é a rejeição editorial precoce (\textit{desk rejection}): a recusa de um manuscrito antes da revisão de mérito, motivada com frequência não por falhas científicas, mas pela inadequação do texto às convenções tácitas de forma e conteúdo esperadas por uma trilha ou conferência específica. Investiga-se, para isso, o projeto de um artefato de software capaz de atuar como tutor de escrita científica. A proposta reposiciona a revisão automatizada: em vez de avaliar o manuscrito contra critérios universais de qualidade, o artefato primeiro modela o gênero comunicativo de uma trilha-alvo a partir do corpus de artigos já aceitos e então verifica a conformidade do manuscrito submetido a essa especificação, gerando diagnósticos acionáveis.
A fundamentação teórica articula a Teoria de Gêneros aplicada a Sistemas de Informação e a análise de \textit{moves} retóricos, na formulação do modelo CARS, integradas sob o método Design Science Research (DSR), envolvendo revisão de literatura, desenvolvimento do artefato e avaliação empírica por meio de estudo
de caso com pesquisadores.
Descreve-se ainda um método eficiente para analisar o conjunto de artigos de uma conferência: extração estruturada por artigo em modelo local, seguida de agregação determinística, minimizando o custo computacional.
% O Simpósio Brasileiro de Sistemas de Informação é adotado como caso de aplicação para a construção e a avaliação inicial do artefato, cuja calibração é derivada do corpus da trilha-alvo e, portanto, transferível a outras conferências mediante corpus adequado. 
Espera-se contribuir com um artefato utilizável, um conjunto de princípios de design generalizáveis e evidência empírica sobre a redução da rejeição precoce e o apoio à formação de pesquisadores no contexto brasileiro.

\textbf{Palavras-chave:}~Escrita Científica;Escrita Acadêmica; Análise de Gêneros; \textit{Moves}
Retóricos; Modelos de Linguagem de Grande Escala (LLM); Design Science Research
\newpage


% Este trabalho investiga o uso de arquiteturas multiagentes baseadas em Inteligência Artificial(IA) como instrumento de apoio à avaliação científica no contexto de periódicos brasileiros. O crescimento da produção científica intensificou limitações estruturais do processo de revisão por pares, como a sobrecarga de revisores, a heterogeneidade dos pareceres e a escassez de retorno estruturado a autores em início de carreira. Uma consequência particularmente onerosa para novos pesquisadores é a rejeição editorial precoce (desk rejection): a recusa de um manuscrito antes da revisão de mérito, motivada com frequência não por falhas científicas, mas pela inadequação do texto às convenções tácitas de forma e conteúdo esperadas por uma trilha ou conferência específica. Este trabalho investiga o projeto de um artefato de software, baseado em uma arquitetura multiagente com modelos de linguagem de grande escala (LLMs), capaz de atuar como tutor de escrita científica. A proposta reposiciona a revisão automatizada, com o objetivo de avaliar o manuscrito contra critérios universais de qualidade, o artefato primeiro modela o gênero comunicativo de uma trilha-alvo a partir do corpus de artigos já aceitos e então verifica a conformidade do manuscrito submetido a essa especificação, gerando diagnósticos acionáveis. 
% A fundamentação teórica articula a Teoria de Gêneros aplicada a Sistemas de Informação e a análise de moves retóricos, integradas sob o método Design Science Research (DSR), envolvendo revisão de literatura, desenvolvimento do artefato e avaliação empírica por meio de estudo de caso com pesquisadores.
% Descreve-se ainda um método eficiente para analisar o conjunto de artigos de uma conferência: extração estruturada por artigo em modelo local, seguida de agregação determinística, minimizando o custo computacional. Espera-se contribuir com um artefato utilizável, um conjunto de princípios de design generalizáveis e evidência empírica sobre a redução da rejeição precoce e o apoio à formação de pesquisadores no contexto brasileiro.




% \textbf{Palavras-chave:}~Revisão Científica Automatizada; Agentes de Inteligência Artificial; Teoria de Gêneros; Design Science Research; Comunicação Científica; Modelos de Linguagem de Grande Escala (LLM)
\newpage
%------------------------------------------------------------------------------
% Fim do resumo
%------------------------------------------------------------------------------
